% =====================================================================================
%  SECTION FILE TEMPLATE                        copy to: <sheet>/sections/<section id>.tex
%  -------------------------------------------------------------------------------------
%  One file per section in sheet.json. build.py writes the section heading, side tab and
%  contents entry from sheet.json, so the file starts with its first \cstopic.
%  Every component is shown once below with MADE-UP content from different kinds of course
%  (physics, history, art history, computer science), so you can see how each one adapts.
%  A real topic uses only the components its content needs. Replace everything.
%  Full reference: references/template-reference.md.
% =====================================================================================

% ---- TOPIC: number, title, where it comes from in the student's material (pages they have)
\cstopic{1.1}{A topic title in 3 to 8 words}{L4 pp.2--9; reading ch.3}
% \allowsplit   <- only if this topic is longer than a page or a split is agreed

% ---- TAKEAWAY: first thing in every topic. What it covers, then the implications.
\takeaway{This topic covers what the reader will find here, in one sentence. Then the two or three implications that answer most questions: what follows, when it applies, how it connects to the next topic.}

% ---- PROSE: a lead-in that says what the next box is for. 2 to 3 sentences.
A short lead-in says what the box below is for and when to use it, so the box is never a bare formula.

% ---- KEY BOX with an equation (math, science, economics, CS)
\begin{keybox}{Name of the result}
\keyeq[What each line shows, in plain words: line 1 gives this, line 2 gives that.]{%
y = a\,x + b \\[4pt]
x = \frac{y - b}{a}}
\begin{symbols}
\(y\); \(x\) & what the symbol is, with units; the other symbol\\
\(a\), \(b\) & slope; intercept\\
\end{symbols}
\means{the one-sentence meaning a student should be able to say out loud.}
\nuance{the condition, exception or common confusion.}
\end{keybox}

% ---- KEY BOX with definitions (any subject)
\begin{keybox}{Terms the exam uses}
\defn{Mercantilism}{the view that a nation grows rich by exporting more than it imports and hoarding bullion.}
\defn{Enclosure}{the fencing of common land into private holdings, mostly in 16th to 18th century England.}
\end{keybox}

% ---- FACT CARD: a work, person, event, case or specimen (humanities, law, biology)
\begin{factcard}{A made-up painting: identify it}[figs/placeholder.pdf][0.36]
\fact{Artist, date}{the artist, c. 1505}
\fact{Medium}{oil on poplar panel}
\fact{Recognize by}{the two or three visual clues that identify it on a slide}
\fact{Significance}{why it matters for the course, in one or two lines}
\end{factcard}
% Without an image: \begin{factcard}{Title} \fact{..}{..} ... \end{factcard}

% ---- CODE (computer science)
\begin{codeblock}[language=Python]
def binary_search(xs, target):
    lo, hi = 0, len(xs) - 1
    while lo <= hi:
        mid = (lo + hi) // 2
        if xs[mid] == target:
            return mid
        lo, hi = (mid + 1, hi) if xs[mid] < target else (lo, mid - 1)
    return -1
\end{codeblock}

% ---- CHANGE TABLE: what happens when you change something, and why. Rows end with \\.
\begin{changetable}
double the sample size & standard error falls by about 30\% & it shrinks like \(\frac{1}{\sqrt{n}}\)\\
raise the price & quantity demanded falls & the demand curve slopes down\\
\end{changetable}
% Other headings: \begin{changetable}[If the facts change][Fact][Outcome][Rule] ... (law)

% ---- TABLE BOX: comparisons, timelines, lookups. Write the rules yourself.
\begin{tablebox}{Compare at a glance}{@{}>{\raggedright\arraybackslash}p{0.22\linewidth}XX@{}}
\toprule & \textbf{Option A} & \textbf{Option B}\\\midrule
cost & low & high\\
when to use & small inputs & large inputs\\
\bottomrule
\end{tablebox}

% ---- FIGURE: width as a fraction of the line, and a caption that says what to notice
\sheetfig[0.4]{figs/placeholder.pdf}{What to notice in the picture, in one sentence.}

% ---- EXAMPLE: a one-line worked use of the box above
\example{with \(a=2\), \(b=1\) and \(y=7\): \(x=\frac{7-1}{2}=3\).}

% ---- WORKED EXAMPLE: formula first, then each part plugged in with a label, then the number
\begin{worked}{Worked example: a title that says what is solved}
\sidebyside[0.45]{%
\small The question, restated in full.\par\medskip
\textbf{Plan:} the approach in one line.}{%
\small
\begin{enumerate}[leftmargin=1.4em]
\item Formula: \(x = \frac{y-b}{a}\).
\item Plug in: \(y = 7\) (given), \(b = 1\) (intercept), \(a = 2\) (slope).
\item Answer: \(x = 3\).
\end{enumerate}}
\end{worked}

% ---- TRAP: a common mistake or a question pattern, with a bold title
\begin{trap}{Select-all questions}
Judge every option on its own; one missed true option usually scores zero.
\end{trap}

% ---- TRUE OR FALSE: 2 to 5 short checks; T or F, the statement, then why in a few words
\begin{tfbox}
\tf{T}{A statement that is true.}{Why, in a few words.}
\tf{F}{A statement with the trap students fall for.}{What is actually true.}
\end{tfbox}
